%\section{Generic constructions for frexlets}
\subsection{Frex via Coproducts with Fral}
\subseclabel{frex-via-coproducts}
\begin{wrapfigure}[4]{r}{3cm}
  \vspace{-1\baselineskip}
\includegraphics{diagram-06.mps}
\end{wrapfigure}
The fral and the frex are related: the free extension of $\A$ by $\X$ is
the \emph{coproduct} of $\A$ with the free algebra over $\X$.
%
A \emph{cospan} $\avar{}$ between $\A_1$ and $\A_2$ consists of two
homomorphisms with a shared codomain: $\A_1
\xto{\aLft}$\aSink$\xleftarrow{\aRgt} \A_2$.  A \emph{cospan
homomorphism} $h : \avar \to \bvar$ is a homomorphism $\UH : \aSink
\to \bSink$ preserving the cospan as in the diagram on the right.  The
coproducts $\A_1\oplus\A_2$ of two models $\A_1$ and $\A_2$ is then
the initial \emph{cospan}.  All models have coproducts, but these may
be difficult to represent. However, in cases such as commutative
monoids, the coproduct is straightforward to represent: its carrier is
the cartesian product of the components carriers.

The universal property of the frex $\A{}[\X]$ combines
those of the fral $\FreeName\,\Pres\!\X$ and
its coproduct with $\A$. Consider the following two diagrams:
\[
\hfill\includegraphics{diagram-05.mps}
\qquad
\hfill\includegraphics{diagram-04.mps}
\]
The \Var-arrows and the \Lft-arrows correspond through the universal
property of the fral: \Lft{} is the unique homomorphic extension of
the \Var{}-arrows. The \Embed-homomorphisms are exactly
the \Rgt-homomorphism.
This identification lets us construct:
\MyCodeListing{
\FrexByCoprodPartA{}\IdrisKeyword{ :}\FrexByCoprodPartB{}\\
\FrexByCoprodPartC{}
}
For commutative monoids it gives the
commutative monoid of linear polynomials with natural numbers as degree-1
coefficients whose carrier is represented by $\CMonoidFrexCarrier{}$.

\subsection{Fral via a Frex}
\subseclabel{fral via frex}
We use the relationship between the fral and frex, described in \subsecref{frex-via-coproducts}, to derive a fral from a corresponding frex.
%
In particular, the following calculation shows that every free algebra arises as a free extension of $\FreeName\,\Pres\emptyset$. We use $\uplus$ for the disjoint union, i.e., the coproduct of sets:
\begin{align*}
\FreeName\,\Pres\X
&\isomorphic \FreeName\,\Pres(\X\uplus\emptyset) & (\X \isomorphic \X \uplus \emptyset)\\
&\isomorphic (\FreeName\,\Pres\X)\oplus(\FreeName\,\Pres\emptyset) & (\text{$\FreeName\,\Pres$ left adjoint, left adjoints preserve coproducts})\\
&\isomorphic (\FreeName\,\Pres\emptyset){}[\X]
\end{align*}
%
Therefore, we may construct a fral from an initial
algebra and its frex:
\MyCodeListing{\PutListing{.5\textwidth}\ByFrexFull{}}
This generic construction can produce suboptimal
representations.  For example, the initial monoid is easy to
construct: its carrier is the unit type. Freely extending this
initial monoid produces alternating lists, that interleave the unit
value. Taking lists of variables instead leads to a simpler representation but
requires more complicated proofs.  However, generic constructions such as the foregoing allow us to trade efficiency for rapid prototyping.

\subsection{Reusing Frexlets}
\subseclabel{reusing frexlets}
The final example demonstrates reuse of one simplifier when
constructing another. Recall the presentation of involutive monoids
from the end of \secref{frexlib}.
\begin{proposition}[Jacobs]\label{prop:involution}
  The carrier set and monoid operations of the free involutive monoid on \X{}
  are those of the free monoid on the product\/ $\BoolX$.
  Similarly, the frex of an involutive monoid by $\X$ is the
  frex of its underlying monoid by\/ $\BoolX$, extended with an involution operation.
\end{proposition}
We can prove this proposition directly, establishing the involutive
axioms, and have taken this strategy in \Frexlib{}.  However, it is
possible to prove this result without referring to the specific
representation of the monoid fral/frex, resulting in a
simplifier-transformer reusing the code implementing simplifiers for
monoids to implement simplifiers for involutive monoid.  The potential
for this kind of reuse extends beyond monoids.  We can phrase the
result in much greater generality, and give a higher-level proof,
using \citepos{jacobs:involution} axiomatisation of involutions. This
more abstract proof generalises to other notions of involutive
algebras, and we plan to exploit it in the future for generic frexlet
reuse. The more abstract proof goes beyond the scope of this
article, involving more abstract category theoretic notions.
